\section{Identificação de Sistemas Instáveis}

A identificação de sistemas é a área voltada à obtenção de modelos a partir de dados experimentais de entrada e saída, conforme apresenta a Figura \ref{identification-model}, em contraposição à modelagem caixa branca, derivada unicamente das leis físicas do processo. Existem dois tipos de identificação de sistemas, o método de caixa preta, onde o modelo é obtido exclusivamente a partir dos dados, e de caixa cinza, que combina uma estrutura física conhecida com a estimação de parâmetros baseada em medições \cite{Aguirre2015}.

\begin{figure}[ht]
    \centering
    \begin{tikzpicture}[scale=1.12, transform shape,
        blk/.style={rectangle, draw=black!70, semithick, rounded corners=2pt, top color=white, bottom color=black!12, minimum height=1.1cm, minimum width=4.6cm, align=center, font=\small},
        arr/.style={-{Stealth[length=2.2mm]}, semithick},
        io/.style={font=\small, align=center},
        ttl/.style={font=\small\bfseries, anchor=west}
    ]
    % ------- a) processo real -------
    \node[ttl] at (-2.2,1.15) {a) Processo em investigação};
    \node[blk, font=\small\bfseries] (sys) at (3.2,0) {Sistema Dinâmico};
    \node[io] (u1) at (-1.2,0) {Entradas\\$u$};
    \node[io] (z1) at (7.6,0)  {Saídas\\$z$};
    \draw[arr] (u1) -- (sys);
    \draw[arr] (sys) -- (z1);
    % ------- b) modelo -------
    \node[ttl] at (-2.2,-1.75) {b) Representação do modelo no domínio do tempo};
    \node[blk, minimum height=2.0cm, bottom color=blue!12] (mod) at (3.2,-3.2)
        {{\bfseries\color{blue!55!black}Modelo Matemático}\\[2pt]
         $\dot{x}(t) = f\big(x(t),u(t),\Theta\big)$\\[1pt]
         $y(t) = g\big(x(t),u(t),\Theta\big)$};
    \node[io] (u2) at (-1.2,-3.2) {$u$};
    \node[io] (y2) at (7.6,-3.2)  {$y$};
    \draw[arr] (u2) -- (mod);
    \draw[arr] (mod) -- (y2);
    \end{tikzpicture}
    \caption{Modelagem Matemática de um Sistema Físico \cite[Adaptado de][]{Jategaonkar2015}.}
    \label{identification-model}
\end{figure}
 

{\color{blue}Formalmente, o problema de identificação parte de uma estrutura de modelo em espaço de estados
\begin{equation}
\dot{\mathbf{x}} = f(\mathbf{x},\mathbf{u},\Theta), \qquad
\mathbf{y} = g(\mathbf{x},\mathbf{u},\Theta), \qquad
\mathbf{z}(t_k) = \mathbf{y}(t_k) + \mathbf{v}(t_k),
\label{eq:prob-ident}
\end{equation}
em que $\mathbf{x}$ é o estado, $\mathbf{u}$ a entrada, $\Theta$ o vetor de parâmetros desconhecidos, $\mathbf{y}$ a saída do modelo e $\mathbf{z}$ a medição, corrompida pelo ruído $\mathbf{v}$. Identificar o sistema é definir a estrutura $f$ e $g$ e estimar $\Theta$ de modo que a saída do modelo reproduza as medições \cite{Jategaonkar2015}.

Essas duas vias de identificação e a modelagem de caixa branca formam um espectro contínuo, no qual as abordagens se distinguem pela origem de cada ingrediente da Equação~\eqref{eq:prob-ident}, conforme ilustra a Figura~\ref{fig:caixas}. A caixa branca, situada no extremo em que nenhum dado de voo é utilizado, constitui modelagem e não identificação, pois estrutura e parâmetros provêm integralmente das leis físicas e de medições diretas, como a geometria, a massa e os ensaios de bancada, procedimento usual em multirrotores convencionais \cite{Bouabdallah2007, Carino2015}, cuja fidelidade fica limitada pelos efeitos que a física simplificada não captura. No extremo oposto, na caixa preta, estrutura e parâmetros são extraídos exclusivamente dos dados, por famílias como os modelos autorregressivos, os métodos de subespaço e as redes neurais. A expressão mais recente dessa classe são os modelos híbridos com aprendizado de máquina, como o NeuroBEM de \textcite{Bauersfeld2021}, que alcança erros de predição muito baixos em voos agressivos de quadrirrotor ao aprender os resíduos aerodinâmicos que a teoria não captura. Essa fidelidade tem um custo estrutural, pois os parâmetros aprendidos não possuem significado físico, não portam limites de confiança e não se conectam de forma direta ao projeto de controladores por alocação de polos, além de exigirem um volume de dados muito superior ao de uma campanha de identificação convencional. A caixa cinza ocupa a posição intermediária, com a estrutura fixada pela física e apenas $\Theta$ estimado dos dados, o que preserva a interpretabilidade e a capacidade de extrapolação do modelo físico e agrega a aderência aos dados da via experimental \cite{Aguirre2015}.

\begin{figure}[H]
  \centering
  \begin{tikzpicture}[
      cx/.style={rectangle, draw=black!70, semithick, rounded corners=2pt, top color=white, bottom color=black!12, minimum height=1.9cm, text width=3.6cm, align=center, font=\small},
      nota/.style={font=\footnotesize, align=center, text=black!75},
      arr/.style={{Stealth[length=2.2mm]}-{Stealth[length=2.2mm]}, semithick}
  ]
  \node[cx] (br) at (0,0) {\textbf{Caixa branca}\\[2pt]\footnotesize estrutura e parâmetros\\ \footnotesize das leis físicas};
  \node[cx, bottom color=blue!12, draw=blue!55!black] (cz) at (4.6,0) {\textbf{Caixa cinza}\\[2pt]\footnotesize estrutura física com\\ \footnotesize parâmetros estimados\\ \footnotesize dos dados};
  \node[cx] (pr) at (9.2,0) {\textbf{Caixa preta}\\[2pt]\footnotesize estrutura e parâmetros\\ \footnotesize extraídos dos dados};
  \node[nota] at (0,-1.55) {geometria, massa,\\ ensaios de bancada};
  \node[nota, text=blue!55!black] at (4.6,-1.55) {\textbf{este trabalho}};
  \node[nota] at (9.2,-1.55) {ARX, subespaço,\\ redes neurais};
  \draw[arr] (-1.9,1.6) -- node[nota, above] {conhecimento físico} node[nota, below] {dependência de dados} (11.1,1.6);
  \end{tikzpicture}
  \caption{\color{blue}Classes de modelo segundo a origem da estrutura e dos parâmetros.}
  \label{fig:caixas}
\end{figure}

}

Este trabalho enquadra-se na abordagem de caixa cinza, pois parte da estrutura física do modelo de corpo rígido, deduzida no Capítulo~\ref{chap:modelagem}, e estima seus parâmetros a partir de dados reais de voo.


O processo de identificação compara, a partir de um mesmo sinal de entrada $u$, a medição do sistema dinâmico e a saída do modelo matemático, para estimação dos melhores parâmetros e estruturas para representar esse processo físico \cite{Aguirre2015}. Porém, como apresentado em \textcite{Jategaonkar2015}, determinados tipos de sistemas são instáveis a entrada em malha aberta, portanto, necessitam de um controlador para que seja possível observar o sistema e realizar a identificação.

Considerando a Figura \ref{identification-diagram} para identificar o sistema $\mathbf{H}$, com o sinal de entrada $u$ e a medição dos sensores $z$, existem duas abordagens diferentes, onde a primeira observa o conjunto completo indicada como identificação por malha-fechada, relacionando as entradas de comando $\boldsymbol{\sigma_P}$ com a saída $z$, a qual é recomendada para processos em que se deseja apenas reproduzir o comportamento global, de comando para resposta. Por reunir a planta e o controlador em um único modelo equivalente, o sistema observado permanece estável, de forma que qualquer método padrão de estimação, como o método do erro de saída, pode ser aplicado sem dificuldades numéricas. Sua principal limitação, contudo, é fornecer apenas esse modelo equivalente do conjunto, sem separar os parâmetros da aeronave da dinâmica do controlador \cite{Jategaonkar2015}.

\begin{figure}[H]
    \centering
    \resizebox{0.78\textwidth}{!}{%
    \begin{tikzpicture}[
        blk/.style={rectangle, draw=black!70, semithick, rounded corners=2pt, top color=white, bottom color=black!12, minimum height=1.3cm, minimum width=2.7cm, align=center, font=\small},
        sum/.style={circle, draw=black!70, semithick, top color=white, bottom color=black!12, minimum size=0.6cm, font=\normalsize},
        arr/.style={-{Stealth[length=2.2mm]}, semithick},
        dbl/.style={{Stealth[length=2.2mm]}-{Stealth[length=2.2mm]}, semithick},
        io/.style={font=\small, align=center},
        lbl/.style={font=\footnotesize, inner sep=2pt},
        dot/.style={circle, fill, inner sep=0pt, minimum size=3.2pt}
    ]
    % ------- malha -------
    \node[io] (cmd) at (-0.85,0.3) {Comando\\$\delta_p$};
    \node[blk] (ctrl) at (2.6,0) {Controlador\\$\mathbf{G}$};
    \node[blk] (air)  at (8.0,0) {Aeronave\\$\mathbf{H}$};
    \node[sum] (ns) at (10.8,0) {$\Sigma$};
    \node[io, text=red!65!black] (ru) at (10.8,1.3) {ruído $v$};
    \node[dot] (tz) at (11.9,0) {};
    \node[io] (z) at (13.3,0) {Resposta\\$z$};
    \draw[arr] (cmd) -- ([yshift=0.3cm]ctrl.west);
    \draw[arr] (ctrl) -- node[lbl, below, pos=0.35] {$u$} (air);
    \draw[arr] (air) -- node[lbl, above, pos=0.45] {$y$} node[lbl, below, pos=0.85] {$+$} (ns);
    \draw[arr] (ru) -- node[lbl, right, pos=0.65] {$+$} (ns);
    \draw[arr] (ns) -- (z);
    \draw[arr] (tz) -- (11.9,-1.6) -- (1.0,-1.6) -- (1.0,-0.3) -- ([yshift=-0.3cm]ctrl.west);
    % ------- faixas de identificacao -------
    \draw[dashed, semithick, black!60] (-0.1,0.55) -- (-0.1,3.15);
    \draw[dashed, semithick, black!60] (11.9,0.55) -- (11.9,3.15);
    \draw[dashed, semithick, black!60] (4.35,0.55) -- (4.35,2.25);
    \draw[dbl] (-0.1,2.9) -- node[fill=white, font=\footnotesize] {Identificação em malha fechada} (11.9,2.9);
    \draw[dbl] (4.35,2.0) -- node[fill=white, font=\footnotesize] {Identificação em malha aberta ($\hat{\mathbf{H}}$)} (11.9,2.0);
    \end{tikzpicture}%
    }
    \caption{Identificação de Sistemas Instáveis com Controlador \cite[Adaptado de][]{Jategaonkar2015}.}
    \label{identification-diagram}
\end{figure}

A segunda abordagem é a identificação em malha aberta dentro do sistema controlado, o qual visa caracterizar diretamente a planta $\mathbf{H}$, o VTOL instável isolado, relacionando a entrada efetiva $u$, aplicada pelo controlador, com a saída medida $z$, como utilizado nos trabalhos de \textcite{Matt2025, Niermeyer2015}. É esta a abordagem adotada neste trabalho, por permitir estimar os parâmetros físicos do próprio drone, e não um modelo equivalente do conjunto planta-controlador. Em contrapartida, a integração de um sistema inerentemente instável pode conduzir à divergência numérica, ou correlação forte entre entrada e saída, quando é empregado o método do erro de saída clássico, exigindo cuidados específicos no processo de estimação \cite{Jategaonkar2015,Tischler2012}.

{\color{blue}
\section[{\color{blue}Métodos de Estimação de Parâmetros}]{Métodos de Estimação de Parâmetros}
\label{sec:estado-arte-metodos}

Na abordagem de caixa cinza, o estimador é o elemento que liga a estrutura física do modelo aos dados experimentais, ao determinar o vetor de parâmetros $\Theta$ que torna a saída simulada compatível com as medições. Esta seção apresenta as principais formas de estimação empregadas na prática aeronáutica e conclui com a justificativa das escolhas desta dissertação. Os métodos aqui formulados operam sobre a estrutura de modelo deduzida no Capítulo~\ref{chap:modelagem} e constituem a base do processo de identificação executado no Capítulo~\ref{chap:identificacao}.

\subsection[{\color{blue}Estimação no domínio da frequência}]{{\color{blue}Estimação no domínio da frequência}}
\label{sec:met-frequencia}

Uma das vias mais empregadas na identificação de aeronaves sustentadas por rotores trabalha no domínio da frequência, consolidada por \textcite{Tischler2012} na metodologia CIFER. O ajuste não é feito sobre o histórico temporal, mas sobre a resposta em frequência estimada a partir de varreduras de excitação, e a coerência entre entrada e saída serve de métrica de qualidade banda a banda, o que permite tratar naturalmente plantas instáveis operando em malha fechada. É a via dominante na identificação de multirrotores de pequeno porte, empregada por \textcite{Wei2017} para levantar o modelo de voo pairado de um quadrirrotor e por \textcite{Ivler2019} para estabelecer diretrizes de ensaio e regras de escalonamento dinâmico entre veículos de portes distintos. O ajuste entrega um conjunto de derivadas de estabilidade e de controle em torno do ponto de operação, entre as quais figuram os coeficientes de amortecimento angular $L_p$, $M_q$ e $N_r$, respectivamente de rolagem, arfagem e guinada, que \textcite{Malpica2024} tratam como grandezas de projeto por condicionarem a agilidade e as qualidades de voo da aeronave.

Em contrapartida, a abordagem impõe exigências de ensaio nem sempre simples de atender. As varreduras precisam ser longas e repetidas em torno de um mesmo ponto de operação, o que demanda espaço de voo amplo e permanência prolongada na condição de interesse, e a resolução nas bandas mais altas depende de telemetria com taxa de amostragem superior à dos registradores embarcados usuais.

\subsection[{\color{blue}Estimação no domínio do tempo}]{{\color{blue}Estimação no domínio do tempo}}
\label{sec:met-tempo}

Outra abordagem forma o resíduo diretamente sobre os históricos temporais medidos, sem passar pela estimativa da resposta em frequência. As exigências de ensaio são, por isso, mais modestas, uma vez que bastam registros sincronizados de entrada e saída ao longo de manobras curtas, compatíveis com campanhas em área confinada e com a taxa de amostragem dos registradores embarcados. Em troca, o esforço desloca-se para a formulação, pois o resíduo é avaliado contra uma estrutura de modelo escolhida de antemão. Dois estimadores organizam essa via e distinguem-se pelo lugar em que o resíduo é formado, o método do erro de equação, que o avalia algebricamente sobre as próprias medições, e o método do erro de saída, que o avalia sobre a saída do modelo integrado no tempo. São apresentados nessa ordem por refletir o seu uso combinado na prática.

\subsubsection[{\color{blue}Método do erro de equação}]{{\color{blue}Método do erro de equação}}

O método do erro de equação, \gls{EEM}, avalia a equação dinâmica diretamente sobre os estados medidos, sem integrar o modelo. O resíduo é a diferença entre a derivada do estado, obtida por diferenciação numérica das medições, e a derivada prevista pela estrutura,
\begin{equation}
\mathbf{e}_{E}(t_k) = \dot{\mathbf{x}}_{m}(t_k) - f\big(\mathbf{x}_{m}(t_k),\,\mathbf{u}(t_k),\,\Theta\big),
\label{eq:eem-geral}
\end{equation}
em que $\mathbf{x}_m(t_k)$ é o estado medido no instante de amostragem $t_k$, $\dot{\mathbf{x}}_m(t_k)$ é a sua derivada, obtida por diferenciação numérica do sinal medido, $\mathbf{u}(t_k)$ é a entrada aplicada ao sistema, $f$ é a estrutura do modelo definida na Equação~\eqref{eq:prob-ident} e $\Theta$ é o vetor de parâmetros a estimar. A Figura~\ref{fig:residuo-eem} ilustra essa formação do resíduo, na qual as medições alimentam diretamente a estrutura, sem nenhuma integração.

\begin{figure}[ht]
  \centering
  \begin{tikzpicture}[
      blk/.style={rectangle, draw=black!70, semithick, rounded corners=2pt, top color=white, bottom color=black!12, minimum height=0.85cm, text width=3.1cm, align=center, font=\footnotesize},
      sum/.style={circle, draw=black!70, semithick, top color=white, bottom color=black!12, minimum size=0.55cm, font=\small},
      arr/.style={-{Stealth[length=2.2mm]}, semithick},
      io/.style={font=\footnotesize, align=center},
      lbl/.style={font=\footnotesize, inner sep=2pt}
  ]
  \node[io] (xm) at (-1.5,0) {$\mathbf{x}_m,\ \mathbf{u}$};
  \node[blk] (fe) at (1.6,0) {estrutura\\ $f(\mathbf{x}_m,\mathbf{u},\Theta)$};
  \node[sum] (se) at (4.6,0) {$\Sigma$};
  \node[io] (dxm) at (4.6,1.2) {$\dot{\mathbf{x}}_m$ (derivada medida)};
  \node[io] (ee) at (7.1,0) {$\mathbf{e}_E$};
  \draw[arr] (xm) -- (fe);
  \draw[arr] (fe) -- node[lbl, below, pos=0.8] {$-$} (se);
  \draw[arr] (dxm) -- node[lbl, right, pos=0.75] {$+$} (se);
  \draw[arr] (se) -- (ee);
  \end{tikzpicture}
  \caption{\color{blue}Formação do resíduo no método do erro de equação, sem integração do modelo.}
  \label{fig:residuo-eem}
\end{figure}

O estimador minimiza a soma dos quadrados desse resíduo. Quando a estrutura é linear nos parâmetros, como ocorre nos modelos de corpo rígido com as inércias fixadas, escreve-se $f = \Phi(\mathbf{x}_m,\mathbf{u})\,\Theta$ e a minimização tem solução fechada de mínimos quadrados \cite{Aguirre2015}
\begin{equation}
\hat{\Theta}_{E} = \Big(\textstyle\sum_k \Phi_k^{\top}\Phi_k\Big)^{-1} \textstyle\sum_k \Phi_k^{\top}\,\dot{\mathbf{x}}_{m}(t_k),
\label{eq:eem-ls}
\end{equation}
em que $\Phi_k = \Phi\big(\mathbf{x}_m(t_k),\mathbf{u}(t_k)\big)$ é a matriz de regressores avaliada no instante $t_k$, que reúne as grandezas conhecidas que multiplicam cada parâmetro, e a soma percorre todas as amostras da janela de dados. As virtudes do método decorrem dessa forma algébrica. O custo computacional é mínimo, não há integração numérica e, por isso, a instabilidade da planta não afeta a estimação. A limitação é estatística, pois os regressores $\Phi$ contêm as próprias medições e o seu ruído, configurando um problema de erros nas variáveis cuja consequência é o viés das estimativas, agravado pela diferenciação numérica, que amplifica o ruído de $\dot{\mathbf{x}}_m$ \cite{Jategaonkar2015}. O EEM é, portanto, o estimador natural de partida, rápido e robusto, mas não o de chegada, papel que cumpre nos trabalhos de \textcite{Graesdal2021} e \textcite{Sato2021}, que o empregam como inicialização do refinamento posterior.

\subsubsection[{\color{blue}Método do erro de saída}]{{\color{blue}Método do erro de saída}}

O método do erro de saída, \gls{OEM}, elimina esse viés ao mudar o lugar onde o resíduo é formado. O modelo é integrado no tempo com a mesma entrada aplicada ao sistema real, produzindo a saída simulada $\mathbf{y}(t_k,\Theta)$, e o resíduo passa a ser o erro de saída $\mathbf{z}(t_k) - \mathbf{y}(t_k,\Theta)$, no qual o ruído de medição aparece apenas na comparação e não dentro dos regressores, como ilustra a Figura~\ref{fig:residuo-oem}, em contraste com a formação algébrica da Figura~\ref{fig:residuo-eem}.

\begin{figure}[ht]
  \centering
  \begin{tikzpicture}[
      blk/.style={rectangle, draw=black!70, semithick, rounded corners=2pt, top color=white, bottom color=black!12, minimum height=0.85cm, text width=3.1cm, align=center, font=\footnotesize},
      sum/.style={circle, draw=black!70, semithick, top color=white, bottom color=black!12, minimum size=0.55cm, font=\small},
      arr/.style={-{Stealth[length=2.2mm]}, semithick},
      io/.style={font=\footnotesize, align=center},
      lbl/.style={font=\footnotesize, inner sep=2pt}
  ]
  \node[io] (u2) at (-1.5,0) {$\mathbf{u}$};
  \node[blk] (mo) at (1.6,0) {modelo integrado\\ $\mathbf{y}(t,\Theta)$};
  \node[sum] (so) at (4.6,0) {$\Sigma$};
  \node[io] (z2) at (4.6,1.2) {$\mathbf{z}$ (medição)};
  \node[io] (eo) at (7.1,0) {$\mathbf{e}_O$};
  \draw[arr] (u2) -- (mo);
  \draw[arr] (mo) -- node[lbl, below, pos=0.8] {$-$} (so);
  \draw[arr] (z2) -- node[lbl, right, pos=0.75] {$+$} (so);
  \draw[arr] (so) -- (eo);
  \end{tikzpicture}
  \caption{\color{blue}Formação do resíduo no método do erro de saída, sobre a saída do modelo integrado.}
  \label{fig:residuo-oem}
\end{figure}



Sob a hipótese de ruído gaussiano branco de covariância $R$, a estimativa de máxima verossimilhança minimiza a log-verossimilhança negativa
\begin{equation}
J(\Theta) = \tfrac{1}{2}\sum_{k=1}^{N} \big[\mathbf{z}(t_k) - \mathbf{y}(t_k,\Theta)\big]^{\top} R^{-1} \big[\mathbf{z}(t_k) - \mathbf{y}(t_k,\Theta)\big],
\label{eq:oem-geral}
\end{equation}
em que $N$ é o número de amostras da janela de dados, $\mathbf{z}(t_k)$ é a medição no instante $t_k$, $\mathbf{y}(t_k,\Theta)$ é a saída simulada pela integração do modelo com os parâmetros correntes e $R$ é a matriz de covariância do ruído de medição, que pondera cada canal pelo inverso da sua variância, dando mais peso às medidas mais confiáveis. O estimador herda as propriedades da máxima verossimilhança, sendo assintoticamente não enviesado e eficiente, e a curvatura do custo fornece, pela matriz de informação de Fisher, o limite de Cramér-Rao para a covariância das estimativas \cite{Jategaonkar2015, Klein2006}. Essas propriedades fazem do OEM o padrão da identificação aeronáutica, aplicado do transporte tripulado aos VANTs de pequeno porte \cite{Tischler2012, Niermeyer2015, Matt2025}. A aplicação a aeronaves de asa rotativa também é consolidada, apesar da maior dificuldade imposta pelo acoplamento entre eixos e pela baixa qualidade das medidas de escoamento em baixa velocidade. \textcite{Jategaonkar2015} dedica uma seção a essa classe, com aplicações dos métodos no domínio do tempo aos helicópteros BO-105 e EC-135, e \textcite{Tischler2012} reconhece o centro aeroespacial alemão como referência mundial na identificação de helicópteros por máxima verossimilhança com sinais do tipo 3-2-1-1. Esses ganhos têm dois custos. A minimização é iterativa e exige integrações repetidas do modelo, e, em plantas instáveis, a integração em malha aberta pode divergir antes do fim da janela de dados, o que demanda estratégias específicas de inicialização e janelamento, apresentadas no Capítulo~\ref{chap:identificacao}. A Figura~\ref{fig:oem-loop} resume o ciclo iterativo do método, no qual o modelo é integrado com a mesma entrada do sistema real, o erro de saída é formado contra a medição e os parâmetros são atualizados pela otimização da função de verossimilhança até a convergência.

\begin{figure}[ht]
    \centering
    \resizebox{0.92\textwidth}{!}{%
    \begin{tikzpicture}[
        blk/.style={rectangle, draw=black!70, semithick, rounded corners=2pt, top color=white, bottom color=black!12, minimum height=1.0cm, minimum width=3.0cm, align=center, font=\small},
        sum/.style={circle, draw=black!70, semithick, top color=white, bottom color=black!12, minimum size=0.6cm, font=\normalsize},
        arr/.style={-{Stealth[length=2.2mm]}, semithick},
        io/.style={font=\small},
        lbl/.style={font=\footnotesize, inner sep=2pt},
        dot/.style={circle, fill, inner sep=0pt, minimum size=3.2pt}
    ]
    % ------- mundo real -------
    \node[io] (ent) at (-0.7,0) {Entrada};
    \node[dot] (d) at (0.5,0) {};
    \node[blk, minimum width=6.8cm, font=\small\bfseries] (sys) at (4.85,0) {Sistema Dinâmico};
    \node[sum] (ns) at (9.6,0) {$\Sigma$};
    \node[io, text=red!65!black] (ru) at (9.6,1.45) {Ruído de medição};
    \node[dot] (tz) at (11.0,0) {};
    \node[io, align=center] (z) at (13.55,0) {Resposta medida\\$z$};
    \draw[arr] (ent) -- (sys);
    \draw[arr] (sys) -- node[lbl, above, pos=0.75] {$+$} (ns);
    \draw[arr] (ru) -- node[lbl, right, pos=0.6] {$+$} (ns);
    \draw[arr] (ns) -- (z);
    % ------- modelo -------
    \node[blk, minimum width=6.8cm, minimum height=2.15cm, bottom color=blue!12] (mod) at (4.85,-2.9) {};
    \draw[semithick, black!70] ([yshift=0.38cm]mod.west) -- ([yshift=0.38cm]mod.east);
    \node[font=\small\bfseries, text=blue!55!black] at ([yshift=0.72cm]mod.center) {Modelo Matemático};
    \node[font=\small, align=center] at ([yshift=-0.36cm]mod.center)
        {{\color{red!70!black}Integração das equações de estado}\\[1pt] Equações de observação};
    \node[sum] (er) at (11.0,-2.9) {$\Sigma$};
    \draw[arr] (d) |- node[lbl, below, pos=0.75] {} (mod.west);
    \draw[arr] (mod.east) -- node[lbl, above, pos=0.5] {$y$} node[lbl, below, pos=0.88] {$-$} (er);
    \draw[arr] (tz) -- node[lbl, right, pos=0.85] {$+$} (er);
    \node[dot] (tz2) at (11.0,-1.35) {};
    \draw[arr] (tz2) -- (0.95,-1.35) -- (0.95,-2.18) -- ([yshift=0.72cm]mod.west);
    % ------- otimizacao -------
    \node[blk, minimum width=6.8cm, minimum height=1.5cm, bottom color=orange!25] (opt) at (4.85,-5.6)
        {Atualização dos parâmetros pela\\ otimização da função de verossimilhança};
    \draw[arr] (er.south) |- node[lbl, right, pos=0.22, align=center] {Erro de resposta\\($z-y$)} ([yshift=-0.25cm]opt.east);
    \draw[arr] ([yshift=-0.75cm]mod.east) -- (9.05,-3.65) -- node[lbl, left, pos=0.5] {Sensibilidades} (9.05,-5.35) -- ([yshift=0.25cm]opt.east);
    \draw[arr] (opt.west) -- (0.5,-5.6) -- node[lbl, left, pos=0.5] {$\Theta_{j+1}$} (0.5,-3.65) -- ([yshift=-0.75cm]mod.west);
    \end{tikzpicture}%
    }
    \caption{Ciclo iterativo do método do erro de saída \cite[Adaptado de][]{Jategaonkar2015}.}
    \label{fig:oem-loop}
\end{figure}


Duas extensões completam a família do domínio do tempo. O método do erro de filtro generaliza o OEM ao incluir ruído de processo por meio de um filtro de Kalman interno ao estimador, sendo indicado para ensaios sob turbulência, ao custo de maior complexidade numérica. As formulações recursivas tratam os parâmetros como estados aumentados de um filtro de Kalman estendido ou \textit{unscented} e permitem estimação em linha durante o voo, úteis para adaptação, mas dependentes da sintonia das covariâncias e sem a mesma transparência estatística da estimação em lote \cite{Jategaonkar2015}. Somam-se a elas as formulações compostas, que articulam mais de um estimador no mesmo procedimento, seja encadeando o erro de equação e o erro de saída, seja confrontando o ajuste temporal obtido por um deles com a resposta em frequência extraída dos mesmos dados, verificação independente recomendada por \textcite{Tischler2012}.

Comum a todas essas variantes está o fato de o resíduo ser formado contra a saída de uma estrutura previamente escolhida, de modo que a qualidade do resultado fica limitada pela fidelidade dessa estrutura. É essa constatação que torna a dedução do modelo não linear da aeronave, objeto do Capítulo~\ref{chap:modelagem}, um passo anterior e necessário à execução da identificação.
}
